% ═════════════════════════════════════════════════════════
%  chapters/ch12-pc-architecture.tex
%  Архитектура ПК-приложения: потоки, мьютексы, очереди
% ═════════════════════════════════════════════════════════

\chapter{Архитектура ПК-приложения}
\label{ch:pc-architecture}

\epigraph{<<UI~--- это не просто кнопки. Это \emph{поток}
событий. Если поток останавливается~--- приложение мертво>>.}{}

\section{Введение}
\label{sec:pc-intro}

В \texttt{ch11} мы создали \emph{первое} GTK4-приложение.
В \texttt{ch05}--\texttt{ch10}~--- \emph{прошивку} для
STM32. Теперь~--- \emph{соединим} их. \textbf{ПК-приложение},
которое \emph{управляет} осциллографом и генератором.

\textbf{Задача}~--- \emph{сложная}. Нужно:
\begin{itemize}
    \item \textbf{Читать} данные из порта \emph{непрерывно}.
    \item \textbf{Рисовать} график \emph{в} \emph{реальном} времени.
    \item \textbf{Реагировать} на \emph{кнопки} пользователя.
    \item \textbf{Не} \emph{зависать} при \emph{этом}.
\end{itemize}

\begin{ideabox}
\textbf{Главная проблема~--- \emph{многопоточность}.} Если \emph{читать}
порт в \emph{главном} потоке~--- \emph{интерфейс} \emph{зависнет}.
Если \emph{рисовать} в \emph{фоновом}~--- \emph{упадёт} (GTK
\emph{не} \emph{потокобезопасен}). Нужен \emph{баланс}.
\end{ideabox}

\subsection{Что вы узнаете}

\begin{itemize}
    \item Почему \emph{один} поток~--- \emph{плохо}.
    \item Как \emph{работают} потоки в GLib.
    \item Что такое \emph{мьютекс} и \emph{зачем} он нужен.
    \item Как \emph{безопасно} обновлять UI из \emph{потока}.
    \item Как \emph{устроен} поток чтения осциллографа.
\end{itemize}

\section{Проблема: один поток}
\label{sec:pc-problem}

\textbf{Представьте}, что мы \emph{пишем} всё в \emph{одном}
потоке:

\begin{listing}[htbp]
\begin{minted}[frame=none, bgcolor=gray!5]{c}
while (1) {
    read(fd, buf, 256);      // блокируется, пока нет данных
    draw_scope(buf);         // рисует график
    gtk_main_iteration();    // обрабатывает события GTK
}
\end{minted}
\caption{Наивный подход: всё в одном потоке}
\label{lst:pc-naive}
\end{listing}

\textbf{Что} \emph{происходит}:
\begin{enumerate}
    \item \texttt{read()} \emph{блокируется}. Пока \emph{нет}
          \emph{данных}~--- программа \emph{висит}.
    \item \emph{Кнопки} \emph{не} \emph{нажимаются}. \emph{Окно}
          \emph{не} \emph{перетаскивается}.
    \item \emph{Пользователь} \emph{думает}, что \emph{программа}
          \emph{сломалась}.
\end{enumerate}

\begin{warnbox}
\textbf{Это} \emph{худший} \emph{вариант}. \emph{Любое}
\emph{блокирующее} \emph{действие} \emph{останавливает}
\emph{весь} UI. \emph{Никакой} \emph{интерактивности}.
\end{warnbox}

\section{Решение: два потока}
\label{sec:pc-solution}

\textbf{Правильный} \emph{подход}~--- \textbf{два} \emph{потока}:
\begin{enumerate}
    \item \textbf{Главный} (UI)~--- \emph{создаёт} \emph{окна},
          \emph{кнопки}, \emph{реагирует} \emph{на} \emph{события}.
    \item \textbf{Фоновый} (чтение)~--- \emph{читает} \emph{порт},
          \emph{складывает} \emph{данные} в \emph{буфер}.
\end{enumerate}


\begin{figure}[htbp]
\centering
\adjustbox{max width=\textwidth}{%
\begin{tikzpicture}[
    x=1cm, y=1cm,
    box/.style={
        draw, thick, rounded corners=1mm,
        minimum width=3cm, minimum height=1cm,
        align=center, font=\footnotesize\ttfamily
    },
    ui/.style={box, fill=clrUIData!20},
    thread/.style={box, fill=clrLogicData!20},
    hw/.style={box, fill=clrHWData!20},
    arr/.style={-{Latex[length=3mm]}, very thick},
    label/.style={
        font=\footnotesize,
        fill=white,
        inner sep=2pt,
        midway,
        above
    },
    labelbelow/.style={
        font=\footnotesize,
        fill=white,
        inner sep=2pt,
        midway,
        below
    }
]

% ─── Координатная сетка (для отладки) ────────────
\draw[step=1cm, gray!20, very thin] (0, 0) grid (11, 5);

% ─── Главный поток (слева) ────────────────────────
\node[ui] (gtk) at (2, 4) {Главный поток\\(GTK)};
\node[ui, below=0.5cm of gtk] (draw) {Рисование\\(Cairo)};

% ─── Фоновый поток (справа) ───────────────────────
\node[thread] (thread) at (8, 4) {Фоновый поток\\(чтение)};
\node[hw, below=0.5cm of thread] (serial) {serial\_port};

% ─── Буфер (посередине, снизу) ────────────────────
\node[thread, minimum width=2cm, minimum height=0.8cm]
    (buf) at (5, 1) {Буфер\\(GMutex)};

% ─── Стрелки внутри потоков ──────────────────────
\draw[arr, color=clrFlowCmd] (gtk) -- (draw);
\draw[arr, color=clrFlowCmd] (thread) -- (serial);

% ─── Стрелки к буферу ────────────────────────────
% write: serial -> buf
\draw[arr, color=clrFlowData] (serial.south west) -- (buf.north east);
% read: buf -> draw
\draw[arr, color=clrFlowData] (buf.north west) -- (draw.south east);

% ─── Подписи (отдельно, не на стрелках) ──────────
\node[font=\footnotesize, above right] at (6.5, 1.5) {write};
\node[font=\footnotesize, above left]  at (3.5, 1.5) {read};

% ─── Координатная сетка (для отладки) ────────────
%\draw[step=1cm, gray!20, very thin] (0, 0) grid (11, 5);
% ─── Метки осей (для отладки) ────────────────────
\foreach \x in {0,...,11} {
    \node[font=\tiny, gray] at (\x, -0.3) {\x};
}
\foreach \y in {0,...,5} {
    \node[font=\tiny, gray] at (-0.3, \y) {\y};
}

\end{tikzpicture}
}% end adjustbox
\caption{Два потока: главный (UI) и фоновый (чтение).
Буфер защищён мьютексом. \emph{Фоновый} поток
\emph{вызывает} \texttt{serial\_port} и \emph{пишет} в буфер.}
\label{fig:pc-two-threads}
\end{figure}






\begin{keybox}
\textbf{Главный поток} \emph{никогда} \emph{не} \emph{блокируется}.
Он \emph{только} \emph{обрабатывает} \emph{события}.
\textbf{Фоновый} \emph{читает} \emph{порт}. \emph{Общаются}~---
\emph{через} \emph{буфер} (**мьютекс**).
\end{keybox}

\section{Мьютекс: защита общих данных}
\label{sec:pc-mutex}

\textbf{Проблема}: \emph{два} \emph{потока} \emph{обращаются} \emph{к}
\emph{одному} \emph{буферу}. \emph{Один} \emph{пишет}, \emph{другой}
\emph{читает}. \emph{Что} \emph{если} \emph{одновременно}?

\begin{itemize}
    \item \emph{Читатель} \emph{возьмёт} \emph{половину} \emph{новых}
          \emph{данных} \emph{и} \emph{половину} \emph{старых}~---
          \emph{график} \emph{<<сломается>>}.
    \item \emph{Или} \emph{читатель} \emph{прочитает} \emph{данные},
          \emph{которые} \emph{ещё} \emph{не} \emph{записаны}
          \emph{до} \emph{конца}.
\end{itemize}

\textbf{Решение}~--- \textbf{мьютекс} (mutual exclusion).
\emph{Замок}, \emph{который} \emph{защищает} \emph{буфер}:
\emph{пока} \emph{один} \emph{поток} \emph{работает}~
--- \emph{другой} \emph{ждёт}.

\begin{listing}[htbp]
\begin{minted}[frame=none, bgcolor=gray!5]{c}
// Инициализация
GMutex mutex;
g_mutex_init(&mutex);

// Писатель (фоновый поток)
g_mutex_lock(&mutex);
// запись в буфер
g_mutex_unlock(&mutex);

// Читатель (главный поток)
g_mutex_lock(&mutex);
// чтение из буфера
g_mutex_unlock(&mutex);
\end{minted}
\caption{\texttt{GMutex}: защита общих данных}
\label{lst:pc-mutex}
\end{listing}

\begin{ideabox}
\textbf{Мьютекс~--- это <<ключ>>.} \emph{Пока} \emph{один}
\emph{поток} \emph{держит} \emph{ключ}~--- \emph{другой}
\emph{ждёт} \emph{у} \emph{двери}. \emph{Когда} \emph{первый}
\emph{отпускает}~--- \emph{второй} \emph{входит}.
\end{ideabox}

\begin{warnbox}
\textbf{Не} \emph{забывайте} \texttt{g\_mutex\_unlock()}. Если
\emph{забыть}~--- \emph{другой} \emph{поток} \emph{зависнет}
\emph{навсегда}. Это \emph{называется} \emph{deadlock}
(\emph{взаимная} \emph{блокировка}).
\end{warnbox}

\section{Поток чтения осциллографа}
\label{sec:pc-reader}

\textbf{Теперь}~--- \emph{конкретно}. \emph{Как} \emph{работает}
\emph{поток}, \emph{который} \emph{читает} \emph{порт}.

\begin{listing}[htbp]
\begin{minted}[fontsize=\footnotesize, frame=none, bgcolor=gray!5]{c}
typedef struct {
    int fd;
    GMutex lock;
    uint8_t buf[8192];
    size_t have;
    gboolean running;
    GThread *thread;
} OscReader;

static gpointer osc_reader_thread(gpointer data)
{
    OscReader *r = data;

    while (r->running) {
        ssize_t n = read(r->fd, r->buf + r->have,
                         sizeof(r->buf) - r->have);
        if (n <= 0) {
            g_usleep(1000);
            continue;
        }
        r->have += n;
        // ... парсинг кадров ...
    }
    return NULL;
}
\end{minted}
\caption{\texttt{osc\_reader\_thread()}: поток чтения}
\label{lst:pc-reader}
\end{listing}

\textbf{Что} \emph{здесь} \emph{важно}:
\begin{itemize}
    \item \texttt{running}~--- \emph{флаг}, \emph{который}
          \emph{останавливает} \emph{поток}.
    \item \texttt{read()} \emph{блокируется}. \emph{Пока} \emph{нет}
          \emph{данных}~--- \emph{поток} \emph{спит}.
    \item \texttt{g\_usleep(1000)}~--- \emph{если} \emph{ошибка}~
          --- \emph{подождать} \emph{1} \emph{мс}, \emph{не}
          \emph{жечь} \emph{CPU}.
    \item \texttt{have}~--- \emph{сколько} \emph{байт} \emph{уже}
          \emph{в} \emph{буфере} \emph{(данные} \emph{могут}
          \emph{приходить} \emph{по} \emph{частям)}.
\end{itemize}

\section{Безопасное обновление UI}
\label{sec:pc-idle}

\textbf{Главная} \emph{проблема}: \emph{UI} \emph{нельзя} \emph{трогать}
\emph{из} \emph{фонового} \emph{потока}. \emph{GTK} \emph{не}
\emph{потокобезопасен}. \emph{Если} \emph{вызвать}
\texttt{gtk\_widget\_queue\_draw()} \emph{из} \emph{потока}~---
\emph{приложение} \emph{упадёт}.

\textbf{Решение}~--- \texttt{g\_idle\_add()}. \emph{Он}
\emph{выполняет} \emph{функцию} \emph{в} \emph{главном} \emph{потоке},
\emph{когда} \emph{он} \emph{свободен}.

\begin{listing}[htbp]
\begin{minted}[frame=none, bgcolor=gray!5]{c}
// Из фонового потока
g_idle_add(on_ui_update, data);

// В главном потоке
static gboolean on_ui_update(gpointer data)
{
    // ... обновляем UI ...
    gtk_widget_queue_draw(app->drawing_area);
    return G_SOURCE_REMOVE;
}
\end{minted}
\caption{\texttt{g\_idle\_add()}: безопасное обновление UI}
\label{lst:pc-idle}
\end{listing}

\begin{keybox}
\textbf{Правило}: \emph{из} \emph{фонового} \emph{потока}~---
\emph{только} \texttt{g\_idle\_add()}. \emph{Всё} \emph{остальное}~
--- \emph{в} \emph{главном} \emph{потоке}.
\end{keybox}

\section{Уроки главы (первая половина)}
\label{sec:pc-lessons-so-far}

\begin{enumerate}
    \item \textbf{Один поток~--- плохо.} UI \emph{зависает}.
    \item \textbf{Два потока~--- хорошо.} UI \emph{отзывчив},
          фоновый \emph{читает}.
    \item \textbf{Мьютекс~--- защита.} Общие данные~---
          \emph{под} \emph{замком}.
    \item \textbf{Не забывайте} \texttt{g\_mutex\_unlock()}.
          Иначе~--- \emph{deadlock}.
    \item \textbf{UI~--- только в главном потоке.}
          Из \emph{фонового}~--- \emph{через} \texttt{g\_idle\_add()}.
\end{enumerate}



%%%%%===============================



\section{Кольцо кадров на ПК}
\label{sec:pc-ring}

В \texttt{ch07} мы сделали \emph{кольцо} \emph{кадров} на
STM32. \emph{Теперь}~--- \emph{то} \emph{же} \emph{самое} на ПК.
\emph{Зачем}? Потому что \emph{данные} \emph{приходят} \emph{быстрее},
чем \emph{рисуются}.

\begin{ideabox}
\textbf{Аналогия: почтовый ящик.} \emph{Курьер} (фоновый поток)
\emph{бросает} \emph{письма} (кадры). \emph{Хозяин} (главный поток)
\emph{читает} \emph{их}, \emph{когда} \emph{свободен}.
\emph{Ящик}~--- \emph{кольцо}. \emph{Если} \emph{полон}~---
\emph{старые} \emph{выбрасываются}.
\end{ideabox}

\subsection{Структура}

\begin{listing}[htbp]
\begin{minted}[fontsize=\footnotesize, frame=none, bgcolor=gray!5]{c}
#define RING_SIZE 16

typedef struct {
    OscFrame frames[RING_SIZE];
    volatile int write_idx;
    volatile int read_idx;
    volatile int dropped;
    GMutex lock;
} FrameRing;
\end{minted}
\caption{Кольцо кадров на ПК}
\label{lst:pc-ring-struct}
\end{listing}

\textbf{Что} \emph{здесь} \emph{важно}:
\begin{itemize}
    \item \texttt{frames[RING\_SIZE]}~--- \emph{массив} \emph{кадров}.
    \item \texttt{write\_idx}~--- \emph{куда} \emph{писать}
          (\emph{фоновый} \emph{поток}).
    \item \texttt{read\_idx}~--- \emph{откуда} \emph{читать}
          (\emph{главный} \emph{поток}).
    \item \texttt{dropped}~--- \emph{счётчик} \emph{потерь}.
    \item \texttt{lock}~--- \emph{мьютекс} \emph{для}
          \emph{защиты}.
\end{itemize}

\subsection{Запись и чтение}

\begin{listing}[htbp]
\begin{minted}[fontsize=\footnotesize, frame=none, bgcolor=gray!5]{c}
// Фоновый поток
bool ring_push(FrameRing *r, const OscFrame *f)
{
    g_mutex_lock(&r->lock);

    int next = (r->write_idx + 1) % RING_SIZE;
    if (next == r->read_idx) {
        r->dropped++;
        g_mutex_unlock(&r->lock);
        return false;
    }
    r->frames[r->write_idx] = *f;
    r->write_idx = next;

    g_mutex_unlock(&r->lock);
    return true;
}

// Главный поток
bool ring_pop(FrameRing *r, OscFrame *out)
{
    g_mutex_lock(&r->lock);

    if (r->read_idx == r->write_idx) {
        g_mutex_unlock(&r->lock);
        return false;
    }
    *out = r->frames[r->read_idx];
    r->read_idx = (r->read_idx + 1) % RING_SIZE;

    g_mutex_unlock(&r->lock);
    return true;
}
\end{minted}
\caption{\texttt{ring\_push()} и \texttt{ring\_pop()}}
\label{lst:pc-ring-push-pop}
\end{listing}

\begin{keybox}
\textbf{Мьютекс~--- обязателен.} \emph{Без} \emph{него}~---
\emph{гонка} \emph{данных}: \emph{читатель} \emph{возьмёт}
\emph{половину} \emph{нового} \emph{кадра} \emph{и}
\emph{половину} \emph{старого}. \emph{График} \emph{<<сломается>>}.
\end{keybox}

\section{Полный поток чтения}
\label{sec:pc-reader-full}

\textbf{Теперь}~--- \emph{полный} \emph{код} \emph{потока}
\emph{чтения}. \emph{Он} \emph{делает} \emph{три} \emph{вещи}:
\emph{читает} \emph{байты}, \emph{парсит} \emph{кадры},
\emph{кладёт} \emph{в} \emph{кольцо}.

\begin{listing}[htbp]
\begin{minted}[fontsize=\footnotesize, frame=none, bgcolor=gray!5]{c}
static gpointer osc_reader_thread(gpointer data)
{
    AppData *app = data;
    uint8_t buf[8192];
    size_t have = 0;

    while (app->thread_running) {
        ssize_t n = read(app->fd_osc, buf + have,
                         sizeof(buf) - have);
        if (n <= 0) {
            g_usleep(1000);
            continue;
        }
        have += n;

        // Парсим все полные кадры
        while (have >= 8) {
            if (!(buf[0] == 0x55 && buf[1] == 0xAA)) {
                memmove(buf, buf + 1, have - 1);
                have--;
                continue;
            }
            uint16_t len = buf[6] | (buf[7] << 8);
            size_t frame_len = 8 + len + 2;
            if (have < frame_len) break;

            // Проверка CRC
            if (crc_ok(buf, frame_len)) {
                OscFrame f;
                parse_frame(buf, &f);
                ring_push(&app->ring, &f);
            }
            memmove(buf, buf + frame_len,
                    have - frame_len);
            have -= frame_len;
        }
    }
    return NULL;
}
\end{minted}
\caption{\texttt{osc\_reader\_thread()}: полный цикл}
\label{lst:pc-reader-full}
\end{listing}

\textbf{Что} \emph{здесь} \emph{важно}:
\begin{itemize}
    \item \textbf{Буфер} \emph{8} \emph{КБ}~--- \emph{достаточно}
          \emph{для} \emph{нескольких} \emph{кадров}.
    \item \textbf{Поиск} \emph{sync}~--- \emph{если} \emph{не}
          \emph{нашли}, \emph{сдвигаем} \emph{на} \emph{1}
          \emph{байт}.
    \item \textbf{Ожидание} \emph{полного} \emph{кадра}~---
          \emph{если} \emph{не} \emph{хватает} \emph{данных},
          \emph{выходим} \emph{из} \emph{цикла}.
    \item \textbf{CRC}~--- \emph{проверяем} \emph{целостность}.
    \item \textbf{Кольцо}~--- \emph{кладём} \emph{готовый}
          \emph{кадр}.
\end{itemize}

\begin{warnbox}
\textbf{Буфер} \emph{может} \emph{переполниться}, \emph{если}
\emph{кадры} \emph{приходят} \emph{быстрее}, \emph{чем}
\emph{обрабатываются}. \emph{Решение}~--- \emph{увеличить}
\emph{буфер} \emph{или} \emph{ускорить} \emph{обработку}.
\end{warnbox}

\section{Остановка потока}
\label{sec:pc-stop}

\textbf{Проблема}: \emph{поток} \emph{работает} \emph{вечно}.
\emph{Как} \emph{его} \emph{остановить} \emph{при} \emph{выходе}?

\textbf{Решение}~--- \texttt{g\_thread\_join()}.

\begin{listing}[htbp]
\begin{minted}[frame=none, bgcolor=gray!5]{c}
// Остановка
app->thread_running = false;   // 1. Флаг
g_thread_join(app->thread);    // 2. Ждём завершения
app->thread = NULL;
\end{minted}
\caption{Остановка потока}
\label{lst:pc-stop}
\end{listing}

\textbf{Что} \emph{происходит}:
\begin{enumerate}
    \item \texttt{thread\_running = false}~--- \emph{поток}
          \emph{увидит} \emph{флаг} \emph{и} \emph{выйдет}
          \emph{из} \emph{цикла}.
    \item \texttt{g\_thread\_join()}~--- \emph{главный}
          \emph{поток} \emph{ждёт}, \emph{пока} \emph{фоновый}
          \emph{завершится}.
\end{enumerate}

\begin{keybox}
\textbf{Без} \texttt{g\_thread\_join()}~--- \emph{утечка}.
\emph{Поток} \emph{может} \emph{обращаться} \emph{к}
\emph{освобождённой} \emph{памяти}. \emph{Приложение}
\emph{упадёт}.
\end{keybox}

\begin{ideabox}
\textbf{Порядок} \emph{важен}: \emph{сначала} \emph{флаг},
\emph{потом} \emph{join}. \emph{Наоборот}~--- \emph{поток}
\emph{никогда} \emph{не} \emph{выйдет} \emph{из} \emph{цикла}.
\end{ideabox}

\section{Уроки главы}
\label{sec:pc-lessons}

\begin{enumerate}
    \item \textbf{Один поток~--- плохо.} UI \emph{зависает}.
    \item \textbf{Два потока~--- хорошо.} UI \emph{отзывчив},
          фоновый \emph{читает}.
    \item \textbf{Мьютекс~--- защита.} Общие данные~---
          \emph{под} \emph{замком}.
    \item \textbf{UI~--- только в главном потоке.}
          Из \emph{фонового}~--- \emph{через} \texttt{g\_idle\_add()}.
    \item \textbf{Кольцо кадров~--- на ПК тоже.}
          \emph{Данные} \emph{приходят} \emph{быстрее}, \emph{чем}
          \emph{рисуются}.
    \item \textbf{Остановка потока~--- через}
          \texttt{g\_thread\_join()}. \emph{Без} \emph{неё}~---
          \emph{утечка}.
\end{enumerate}

\section{Чек-лист грабель}
\label{sec:pc-pitfalls}

\begin{table}[htbp]
\centering
\caption{Грабли многопоточности}
\label{tab:pc-pitfalls}
\small
\begin{tabularx}{\textwidth}{@{}c L L@{}}
\toprule
\textbf{№} & \textbf{Грабля} & \textbf{Решение} \\
\midrule
1 & UI зависает & Чтение в отдельном потоке \\
2 & Гонка данных & \texttt{GMutex} \\
3 & Забыли \texttt{g\_mutex\_unlock()} & Deadlock \\
4 & Обращение к UI из потока & \texttt{g\_idle\_add()} \\
5 & Забыли \texttt{g\_thread\_join()} & Утечка \\
6 & \texttt{thread\_running} не \texttt{volatile} & Поток не видит изменений \\
7 & Кольцо переполнено & \texttt{dropped++} \\
8 & Нет проверки CRC & Мусор в кадре \\
\bottomrule
\end{tabularx}
\end{table}

\section{Упражнения}
\label{sec:pc-exercises}

\begin{exercisebox}
\textbf{Упражнение 1.} Что произойдёт, если \emph{убрать}
\texttt{g\_mutex\_lock()} из \texttt{ring\_push()}? Приведите
сценарий с \emph{гонкой}.
\end{exercisebox}

\begin{exercisebox}
\textbf{Упражнение 2.} Реализуйте \emph{два} фоновых потока:
один читает \emph{осциллограф}, второй~--- \emph{генератор}.
\emph{Как} \emph{разделить} \emph{мьютексы}?
\end{exercisebox}

\needspace{5\baselineskip}
\begin{exercisebox}
\textbf{Упражнение 3.} Добавьте \emph{счётчик} \emph{потерь}
в \emph{статусбар}. Как \emph{обновлять} его \emph{без}
\emph{блокировки} \emph{UI}?
\end{exercisebox}

\begin{exercisebox}
\textbf{Упражнение 4.} Предложите способ \emph{ускорить}
\emph{парсинг} \emph{кадров}, если \emph{данных} \emph{много}.
\emph{Что} \emph{изменить} в \emph{буфере}?
\end{exercisebox}

\section{Проверьте сами}
\label{sec:pc-check}

\begin{enumerate}
    \item Соберите \emph{ПК-приложение}.
    \item Подключите \emph{Nucleo} с \emph{осциллографом}.
    \item Запустите \emph{поток} \emph{чтения}.
    \item \emph{Проверьте}, что \emph{UI} \emph{не} \emph{зависает}.
    \item \emph{Остановите} \emph{поток}. \emph{Убедитесь}, что
          \emph{приложение} \emph{закрывается} \emph{чисто}.
    \item \emph{Сделайте} \emph{скриншот}~--- \emph{пришлите}
          \emph{в} \emph{книгу}.
\end{enumerate}

\begin{ideabox}
\textbf{Если получилось~--- вы поняли многопоточность.}
Это \emph{главное}. Дальше~--- \emph{отрисовка} \emph{графика},
\emph{интерфейс} \emph{генератора}. \emph{Всё} \emph{впереди}.
\end{ideabox}

